\section{Trois lectures ontologiques (R1, R2, R3)}
\label{sec:lectures_ontologiques}

La PT, par sa nature de théorie purement informationnelle dérivée
d'un unique paramètre arithmétique $s = 1/2$, admet trois lectures
ontologiques compatibles avec sa structure formelle. Cette
trichotomie est exposée en détail dans les chapitres~24 et~26 de
la monographie~\cite{PT_Monograph} ; nous la reprenons ici
brièvement, et l'appliquons spécifiquement au débat
Aharonov--Bohm / Bell développé dans cet article. Le point
crucial à souligner --- et qui constitue l'essentiel
de~\S~\ref{ssec:independance} --- est que les démonstrations des
sections~\ref{sec:champ_connexion}--\ref{sec:prediction_qh1a}
sont logiquement indépendantes du choix de lecture. La trichotomie
porte sur l'\emph{interprétation} de ce que les théorèmes
disent, non sur leur validité.

\subsection{R1 : la PT comme coïncidence mathématique}
\label{ssec:R1}

La lecture R1 considère que la PT reproduit le bon contenu
empirique --- la phase Aharonov--Bohm, les corrélations EPR,
l'absence de dépendance géométrique de la MI --- par une
coïncidence structurale entre la combinatoire du crible
d'Eratosthène et la structure phénoménologique de
l'électrodynamique quantique. Aucun engagement ontologique n'est
pris sur la nature physique des angles $\theta_p$ : ils sont des
quantités mathématiques utiles à la prédiction, sans plus. Les
théorèmes des sections précédentes restent valides --- ils sont
des assertions mathématiques sur la structure du crible --- mais
leur correspondance avec les phénomènes observés est traitée comme
un fait empirique inexpliqué.

Cette lecture est la plus prudente épistémologiquement : elle
n'engage que le minimum nécessaire pour rendre compte des
observations. Son point faible est qu'elle ne fournit aucune
\emph{explication} de la coïncidence : pourquoi la combinatoire
du crible donne-t-elle exactement la bonne structure
phénoménologique ? La lecture R1 ne répond pas à cette question.

\subsection{R2 : structuralisme}
\label{ssec:R2}

La lecture R2 considère que le crible d'Eratosthène et
l'électrodynamique quantique partagent un \emph{pattern}
holonomique abstrait, et que ni l'un ni l'autre n'est
ontologiquement premier. Les $\theta_p$ ne sont alors
ni des objets mathématiques (au sens R1) ni des objets physiques
(au sens R3) : ils sont les composantes d'une structure abstraite
dont les deux théories sont des incarnations.

Cette lecture, classique dans la philosophie des sciences sous
le nom de \emph{structuralisme} \cite{Belot1998,Maudlin2011},
trouve en PT un terrain particulièrement favorable : la
correspondance entre le crible et la phénoménologie n'est plus
une coïncidence mais l'expression d'une structure commune
sous-jacente, et la primauté de l'holonomie devient une propriété
structurelle plutôt qu'une question d'ontologie matérielle. Les
théorèmes des sections précédentes conservent leur validité, et
acquièrent un statut explicatif au sens où ils décrivent la
structure commune.

\subsection{R3 : réalisme informationnel}
\label{ssec:R3}

La lecture R3 considère que les $\theta_p$ \emph{sont} la réalité
physique fondamentale, et que les objets classiques de
l'électrodynamique --- $A_\mu$, $F_{\mu\nu}$, $g_{\mu\nu}$ ---
sont des projections phénoménologiques sur l'espace-temps
émergent. Sous cette lecture, la substance ontologique de la
théorie est l'information ($\DKL$ en bits) ; les notions de
particule, champ, espace, et temps sont des constructions
émergentes au-dessus de cette substance informationnelle.

R3 est la lecture la plus engagée : elle affirme que le monde est
fait d'information arithmétique, que la matière et l'espace n'en
sont que des projections, et que les corrélations observées dans
les tests de Bell sont des reflets directs de la structure CRT
fondamentale. Sous R3, la dissolution de la non-localité spatiale
(\S~\ref{ssec:dissolution_consequence}) prend son sens le plus
fort : les corrélations EPR sont littéralement non-spatiales,
parce que ce qui est mesuré n'a pas de support spatial.

R3 hérite des difficultés philosophiques classiques du réalisme
informationnel \cite{Healey2007} : la nature de
l'« exécution » du crible reste ouverte (qui ou quoi
« calcule » la suite des $\{g_n\}$?), et l'apparente
correspondance entre la structure mathématique et la
phénoménologie observée demande, sous R3, une explication
métaphysique qui dépasse la portée de la théorie elle-même.

\subsection{Indépendance des théorèmes par rapport à la lecture}
\label{ssec:independance}

Le point logique central de cette section est que les
Théorèmes~\ref{thm:invariance_geom}
et~\ref{thm:ruelle_zero_fr}, le
Corollaire~\ref{cor:closed_form_fr}, et la
Prédiction~\ref{pred:qh1a} sont \emph{tous} valides sous R1, R2,
R3. Aucun résultat formel des sections
précédentes ne dépend du choix d'interprétation ontologique. Les
démonstrations sont mathématiques ; les prédictions sont des
fonctions des angles d'holonomie et de la métrique émergente ; et
les valeurs numériques (par exemple
$\mathcal{I}_{3,5} = 0{,}138$ bits) sont calculables
indépendamment de la lecture choisie.

Ce qui change avec la lecture, c'est uniquement
l'\emph{interprétation} de ce que mesure l'expérience. Sous R1,
une confirmation expérimentale de QH1a est un fait empirique
inexpliqué de coïncidence avec la structure du crible. Sous R2,
elle est l'expression empirique d'une structure abstraite
commune entre crible et électrodynamique. Sous R3, elle est une
mesure directe de la réalité ontologique informationnelle. Les
trois lectures sont compatibles avec les mêmes observations ; le
choix entre elles dépend de critères extra-empiriques (parcimonie,
pouvoir explicatif, engagement métaphysique).

Cette indépendance a une conséquence pratique importante pour
l'évaluation expérimentale de la PT : un test de QH1a peut être
mené, et ses résultats interprétés, sans engagement préalable sur
le statut ontologique des $\theta_p$. La théorie est testable
indépendamment de la lecture. C'est, à notre connaissance, le
premier cadre théorique qui propose une dissolution structurelle
du débat Aharonov--Bohm / Bell tout en restant
empiriquement falsifiable, et qui peut être évalué sans avoir à
choisir d'abord entre les positions ontologiques classiques.
